% Generated draft: requires mathematical and source review.
\documentclass[11pt]{article}
\usepackage[margin=25mm]{geometry}
\usepackage{fontspec}
\setmainfont{Latin Modern Roman}
\usepackage{amsmath,amssymb,mathtools,unicode-math}
\setmathfont{Latin Modern Math}
\usepackage{enumitem,needspace,xurl,hyperref}
\setlength{\parindent}{0pt}
\setlength{\parskip}{5pt}
\setlength{\emergencystretch}{3em}
\newcommand{\N}{\mathbb{N}}
\newcommand{\Z}{\mathbb{Z}}
\newcommand{\Q}{\mathbb{Q}}
\newcommand{\R}{\mathbb{R}}
\newcommand{\C}{\mathbb{C}}
\newcommand{\F}{\mathbb{F}}
\newcommand{\A}{\mathbb{A}}
\newcommand{\PP}{\mathbb{P}}
\newcommand{\E}{\mathbb{E}}
\newcommand{\eps}{\varepsilon}
\newcommand{\dd}{\,\mathrm{d}}
\begin{document}
\section*{UnsolvedMath ID: 30000826}
\textbf{Source problem code:} catalogue entry
\subsection{Definitions and mathematical statement}
For a path-connected space \(X\), let \(PX=C([0,1],X)\) be its path space and let \(\mathrm{ev}:PX\to X\times X\) send \(\gamma\) to \((\gamma(0),\gamma(1))\). In the unreduced convention, \(\operatorname{TC}(X)\) is the least number of open sets covering \(X\times X\) on each of which \(\mathrm{ev}\) has a continuous section; it is infinite if no finite such cover exists. The record asks two questions.

First, for a locally trivial fibration \(F\to E\to B\), is
\[
\operatorname{TC}(E)\leq\operatorname{TC}(F)\operatorname{TC}(B)?
\]
Local triviality means that the projection is locally isomorphic over each trivializing open subset \(U\subseteq B\) to \(U\times F\to U\). The intended hypotheses on connectedness and on the category of spaces are not supplied. In particular, no product trivialization or restriction on the structure group is part of the extracted question.

Second, for the closed connected nonorientable surface
\[
N_g=\underbrace{\R\PP^2\mathbin{\#}\cdots\mathbin{\#}\R\PP^2}_{g\text{ summands}},\qquad g\geq1,
\]
determine \(\operatorname{TC}(N_g)\) for every genus \(g\).

The record does not state the normalization of topological complexity. If \(\operatorname{TC}_{\mathrm{red}}=\operatorname{TC}-1\), the displayed inequality in the unreduced reading becomes
\[
\operatorname{TC}_{\mathrm{red}}(E)+1
\leq(\operatorname{TC}_{\mathrm{red}}(F)+1)(\operatorname{TC}_{\mathrm{red}}(B)+1).
\]
It must not be silently replaced by the product inequality with reduced values in all three positions. Likewise, all requested surface values depend on this normalization. The closed-surface convention is the customary reading of the supplied genus notation and requires source confirmation if surfaces with boundary were intended.
\subsection{Short English statement}
Does the topological complexity of a locally trivial bundle satisfy the stated product bound in terms of fibre and base? Separately, determine the topological complexity of each closed nonorientable surface. The normalization of \(\operatorname{TC}\) must be fixed.
\subsection{Sources}
\begin{itemize}
\item Local UnsolvedMath registry, permanent ID \(30000826\), source code catalogue entry, ``Topological Complexity of Fibrations and Nonorientable Surfaces.'' No individual proposer is supplied.
\item \emph{Mini-Workshop: Topology of closed one-forms and Cohomology Jumping Loci} (2008), Oberwolfach Reports, DOI 10.4171/owr/2007/40, pp. 2321--2360: \url{https://doi.org/10.4171/owr/2007/40}.
\item Supplied reference, ``nonorientable surfaces'': \url{https://www.sciencedirect.com/science/article/pii/S0166864117304777}.
\item Supplied reference, ``fibre bundle bounds'': \url{https://intlpress.com/api/bgcloud-front/resource/pdf/volume/1844209401469718529-1844209401469718529-9f4d3018c2e31aca2e3887988363e401.pdf}.
\item The literature triage checked 2026-08-21 labels the combined item partially solved, says the fibration inequality has refinements, and reports surface results in several genera with low genera unresolved in the cited classification. It supplies no exact genus list or values; no present-day status is inferred from that historical description.
\end{itemize}
\end{document}
